\section{Related work, and the gap}
\label{sec:related}

Four literatures meet at a demand ceiling, and the gap sits where they fail to
overlap.

\paragraph{Stochastic and chance-constrained MPC for buildings.}
\citet{oldewurtel2012} and \citet{oldewurtel2014} established chance-constrained
MPC for building climate control under weather uncertainty, reformulating
per-step chance constraints deterministically under Gaussian disturbance
assumptions; \citet{zhang2013scenario} replaced the distributional assumption with
sampled scenarios in the sense of \citet{calafiore2006scenario}. Joint-in-time
chance constraints --- that a trajectory stay feasible at \emph{every} step of
the horizon --- are handled in the general stochastic-MPC literature by risk
allocation under Boole's inequality \citep{ono2008risk,paulson2017joint},
which is conservative by construction and, as Section~\ref{sec:horizon}
measures, vacuous at the horizon a demand ceiling requires. A recent line
combines conformal prediction with MPC to obtain joint-in-time regions without
a distributional assumption \citep{lindemann2023safe,vogel2025conformalmpc},
evaluated on trajectory-prediction and synthetic linear-system benchmarks. The
one application of that idea to building demand response we are aware of
\citep{xu2026e2edpc} uses stage-wise --- that is, marginal --- conformal
tightening in an EnergyPlus simulation and does not report horizon-level
violation. None of these measures, on metered building load, how far a
marginally calibrated $q95$ constraint's realised horizon risk sits from its
nominal level.

\paragraph{Demand-charge management.} The closest prior work on monthly peak
limiting under load uncertainty is the Hong Kong PolyU line
\citep{xu2019adaptive,xu2020exploration}: probabilistic ANN load forecasts,
driven by probabilistic weather forecasts, feed an expected-cost threshold
resetting scheme, and Monte Carlo is used to study the robustness of the
\emph{threshold}, not to bound the probability of exceeding it.
\citet{flamm2020peak} treat demand-charge management with storage as
multi-period stochastic MPC over weather scenarios with affine decision rules
and an expected-cost objective, with a weighting for horizons that straddle
billing periods. In neither line is the exceedance probability of the ceiling
a constraint with a stated level, and neither is evaluated on an Indian tariff
or an Indian building. Field and simulation studies of HVAC demand flexibility
under MPC report peak reduction and savings without a risk statement at all.

\paragraph{Conformal prediction for multi-step forecasts.} Split conformal
\citep{vovk2005algorithmic} and CQR \citep{romano2019cqr} give per-step
marginal coverage. \citet{stankeviciute2021conformal} obtain joint coverage
over a horizon by Bonferroni correction; \citet{sun2024copula} replace
Bonferroni with a copula over per-step conformity scores, which is the same
modelling move as Section~\ref{sec:dial} but deployed as a \emph{predictor}
rather than inside an optimiser, and requiring a second calibration set.
\citet{lin2022temporal} study conformal intervals under temporal dependence.
\citet{zaffran2022adaptive} and \citet{wang2024acmcp} extend adaptive conformal
inference \citep{gibbs2021adaptive} to non-stationary and multi-step settings,
targeting marginal coverage; the latter establishes that optimal $h$-step
errors are serially correlated to lag $h-1$, which is the dependence our
copula fits. This paper adds to that line a measurement rather than a method:
the realised joint exceedance on real load errors, placed against both the
Bonferroni bound and the independence figure, with the copula's accuracy for
that quantity reported.

\paragraph{Decision-aware forecasting and calibration.} Decision-focused
learning \citep{elmachtoub2022spo,donti2017task,berthet2020perturbed} trains
the forecast on the downstream decision's cost; \citet{stratigakos2026decision}
calibrate multivariate prediction sets to the \emph{reliability of the
downstream decision} rather than to coverage, for robust reserve scheduling.
That is complementary to Section~\ref{sec:dial}, which reaches the same object
--- an operator-set violation level on a committed ceiling --- from the
optimiser side, and reports its exactness and resolution floor.

\paragraph{Data.} Probabilistic load forecasting is mature at system scale
\citep{hong2016review,hong2016gefcom} and increasingly benchmarked at building
scale \citep{miller2020bdg2,emami2023buildingsbench,nweye2023citylearn}. Aggregation is known to
improve point accuracy by a scaling law \citep{sevlian2018scaling}; whether
\emph{distribution-free coverage} holds across aggregation levels has not, to
our knowledge, been reported. For India, \citet{rashid2019iblend} publish
I-BLEND: 52 months of one-minute mains data, with power factor, for seven
buildings on a Delhi campus. It is the only open Indian building-level load
corpus at the resolution Indian demand charges are assessed on, and it has
been used for point forecasting rather than for demand-charge control.

\paragraph{The gap.} We find no published work that (i) measures, on metered
building load, the horizon-level breach probability a marginally calibrated
$q95$ chance constraint actually delivers against its nominal level, bracketed
by the valid Boole bound above and the independence figure that is usually
misquoted as one; (ii) restores that level as a settable parameter inside a
demand-charge MILP and reports its exactness and resolution floor; (iii) tests
distribution-free coverage across building and system scale under one
protocol; or (iv) does any of this on real Indian commercial buildings at
native 15-minute resolution under an Indian tariff. Items (i)--(iii) are this
paper. Item (iv) is the work Section~\ref{sec:india-fix} opens, and the data
for it exists.
